\documentclass[11pt]{article}
\usepackage{amsmath,amssymb,amsthm,hyperref}
\newtheorem{theorem}{Theorem}
\newtheorem{lemma}[theorem]{Lemma}
\newtheorem{corollary}[theorem]{Corollary}
\title{An Extreme-Root Obstruction to Ideal Permutation Variations}
\author{}
\date{November 23, 2026}
\begin{document}
\maketitle

\paragraph{Scope.}
This note closes the fixed-root, ideal-first-variation construction as
a route to the universal target. It permits arbitrary distinct real
roots, chosen separately for each target, and arbitrary affine
permutation lists. The result is not a lower bound for general head
complexity. In particular, it does not apply when the common root
multiset itself varies with the Boolean input. The proofs below are
manual, not Lean formalizations.

\begin{lemma}[Minimum fibers are faces]
Let $L(x)=a_0+\sum_i a_i x_i$ on $\{0,1\}^n$, and suppose
$L(x)\ge\lambda$ everywhere. Then $\{x:L(x)=\lambda\}$ is
either empty or a coordinate subcube.
\end{lemma}
\begin{proof}
The minimum is $a_0+\sum_{a_i<0}a_i$. If this is larger than
$\lambda$, the fiber is empty. Otherwise it equals $\lambda$, and
\[
 L(x)-\lambda=\sum_{a_i>0}a_i x_i+
                    \sum_{a_i<0}(-a_i)(1-x_i).
\]
Every summand is nonnegative. Their sum vanishes exactly when
$x_i=0$ for $a_i>0$ and $x_i=1$ for $a_i<0$; coordinates with
zero coefficient are free.
\end{proof}

\begin{theorem}[Two face partitions determine every ideal variation]
Let $\lambda_1<\cdots<\lambda_m$ be real numbers. Let
$u_1,\ldots,u_m$ and $v_1,\ldots,v_m$ be affine forms such that
each of their value lists is a permutation of these numbers at every
Boolean input. Let $a_h,b_h$ be affine forms, and suppose
\[
 \left.\partial_\epsilon\right|_0
 \left[\prod_h(s+u_h(x)+\epsilon a_h(x))-
       \prod_h(s+v_h(x)+\epsilon b_h(x))\right]=r(x)
\tag{1}
\]
is independent of $s$. There are two partitions of the cube into
at most $m$ coordinate subcubes such that, on each cell of their
common refinement, $r$ is affine. More precisely, if
$U_h=\{x:u_h(x)=\lambda_1\}$ and
$V_h=\{x:v_h(x)=\lambda_1\}$, then
\[
 r(x)=d\left(\sum_h\mathbf1_{U_h}(x)a_h(x)
                   -\sum_h\mathbf1_{V_h}(x)b_h(x)\right),
 \qquad d=\prod_{j=2}^m(\lambda_j-\lambda_1)>0.
\tag{2}
\]
The empty product for $m=1$ is one.
\end{theorem}
\begin{proof}
Every $u_h$ and $v_h$ is at least $\lambda_1$ on the cube, so
the lemma applies. Exactly one form in each list attains
$\lambda_1$ at each vertex. Thus the nonempty fibers are disjoint
partitions. Evaluate (1) at $s=-\lambda_1$. In each product
derivative, all terms vanish except the term differentiating its
unique zero factor. The remaining product is $d$, independently of
the input. This gives (2). An intersection of two coordinate subcubes
is empty or a coordinate subcube, and on that intersection (2) is a
fixed difference of two affine forms.
\end{proof}

\begin{corollary}[Uniform counting obstruction]
Write $N=2^n$ and $p=2m(n+1)\le N$. Even allowing the roots,
the two lists, and all perturbations to depend on the target, the
number of strict sign tables obtainable as $\operatorname{sgn}r$
in (1) is at most
\[
 (3^n+1)^{2m}\left(\frac{eN}{p}\right)^p.
\tag{3}
\]
Consequently, for every fixed $C>0$, the union over
$m\le C2^n/n^2$ gives only $2^{o(2^n)}$ sign tables as
$n\to\infty$. This is insufficient for the requested universal bound.
\end{corollary}
\begin{proof}
There are $3^n$ nonempty coordinate subcubes, specified by making
each coordinate free, zero, or one. Allowing an empty fiber gives
at most $(3^n+1)^{2m}$ ordered pairs of fiber lists. We may
overcount lists that do not form partitions.

For each pair that does form partitions, (2) is linear in the
$p=2m(n+1)$ affine coefficients. Its strict sign tables are the
regions of an arrangement of $N$ linear hyperplanes in
$\mathbb R^p$. The usual deletion-restriction induction gives at
most $\sum_{j=0}^p\binom Nj$ regions, even for affine hyperplanes:
adding a hyperplane can split no more regions than the induced
arrangement of the preceding hyperplanes on it. This recurrence
and the initial cases prove the stated sum bound.
For $t=p/N\in(0,1]$, positivity of binomial coefficients gives
\[
 \sum_{j=0}^p\binom Nj\le t^{-p}(1+t)^N
                         \le(eN/p)^p.
\]
The positive factor $d$ in (2) has no effect on signs. In particular,
varying the real root values contributes no extra sign patterns once
the face lists have been specified. This proves (3).

At $m\le C N/n^2$, the base-two logarithm of the first factor is
$O_C(N/n)$. The logarithm of the second is
$O_C(N\log n/n)$, uniformly over these $m$, by the monotonicity
of $p\log(eN/p)$ for $0<p\le N$. Summing over the at most
$C N/n^2$ permitted integers $m$ adds only $O_C(n)$ to the
logarithm. The result is $o(N)$.
\end{proof}

\begin{corollary}[Parity forces many roots]
If the output in (1) has the strict signs of $n$-bit parity, or its
complement, then $m^2\ge2^{n-1}$.
\end{corollary}
\begin{proof}
An affine function cannot strictly represent parity on a coordinate
subcube of dimension at least two. Indeed, its signed sum against
the parity character is zero, since that character is orthogonal to
the constant and every coordinate, whereas strict representation
would make every summand positive after possibly reversing a fixed
sign. Therefore every nonempty cell of the common refinement in the
theorem has at most two vertices. There are at most $m^2$ cells,
so $2^n\le2m^2$.
\end{proof}

\begin{theorem}[Higher parameter order at simple roots does not evade the obstruction]
Replace each perturbed factor in (1) by
\[
 s+u_h(x)+\sum_{q\ge1}\epsilon^q a_{h,q}(x),
 \qquad
 s+v_h(x)+\sum_{q\ge1}\epsilon^q b_{h,q}(x),
\]
where every coefficient is affine in $x$. Finite polynomials or
formal power series in $\epsilon$ are both allowed. Suppose all
coefficients of the product difference of orders less than $k$ vanish
as polynomials in $s$ at every Boolean input, and the coefficient of
order $k$ is $r(x)$, independent of $s$. Then (2) holds with
$a_h=a_{h,k}$ and $b_h=b_{h,k}$. In particular (3) and the parity
bound still apply, independently of $k$.
\end{theorem}
\begin{proof}
Fix a Boolean input and match the two initial factor lists by their
distinct root values. We show inductively that the coefficients of
the matched factors agree at every order $q<k$. If agreement holds
at all smaller orders, then the coefficient of order $q$ in the
product difference is
\[
 \sum_h(a_{h,q}(x)-b_{\pi(h),q}(x))
                  \prod_{\ell\ne h}(s+u_\ell(x)).
\]
All terms involving only lower coefficients cancel in matched pairs.
Evaluation at $s=-u_h(x)$ and distinctness of the roots force
$a_{h,q}(x)=b_{\pi(h),q}(x)$. The induction starts at $q=1$.
The same displayed formula consequently holds at order $k$, whether
or not its coefficient vanishes. It is exactly the first variation
in (1), using the affine order-$k$ coefficient forms. Apply the
theorem.
\end{proof}

\paragraph{Repeated roots.}
A nonzero ideal first variation cannot start from a repeated root:
evaluating its product derivative at that root gives zero. Higher
orders at repeated roots are outside the final theorem. So are
input-dependent common root multisets and variations that are not
constant in the auxiliary variable $s$. No universal construction
in any of these larger classes is asserted here.

\paragraph{Research consequence.}
The four-bit weighted example remains a correct nonaffine identity.
Its root lists and perturbations were independently checked exactly.
However, the extreme-root invariant proves that extending that
fixed-root ideal-variation strategy cannot settle the universal target.
The family is therefore closed as a universal approach, rather than
being left behind an unproved existence lemma.
\end{document}
